\documentclass[a4paper,11pt]{article}
\usepackage{../../exercise-style-cmbright}

\renewcommand{\BookTitle}{Precalculus 2e (OpenStax, 2021)}
\renewcommand{\ExerciseSource}{Section}
\renewcommand{\SectionLabel}{1.1}
\renewcommand{\ExerciseNumber}{3}
\renewcommand{\AuthorName}{Henrik Samuelsson}
\renewcommand{\ProblemTitle}{Exercise Solution}

\begin{document}
\maketitle

\begin{problem}
Why does the vertical line test tell us whether the graph of a relation represents a function?
\end{problem}

\begin{solution}
On a graph, the \(x\)-coordinate represents the input and the \(y\)-coordinate represents the output. A vertical line fixes one input value because every point on that line has the same \(x\)-coordinate. The points where the line intersects the graph show the outputs---their \(y\)-coordinates---for that input.

For example, if a graph contains the points \((2,1)\) and \((2,4)\), both points have the input \(x=2\), but they have different outputs, \(y=1\) and \(y=4\). The vertical line \(x=2\) intersects the graph at both points, showing that the input has more than one output.
\end{solution}

\begin{answer}
If a vertical line intersects the graph of a relation more than once, the same input has more than one output. For the relation to be a function, each input must have exactly one output.
\end{answer}

\end{document}
